% Source: https://artofproblemsolving.com/community/c3h1734824_analytic_continuation
% Stable identity: aops-topic-1734824-q01
An infinite series means the limit of its partial sums, when that limit exists.

\textbf{First series.} The partial sums are $S_n=\dfrac{1}{2}+\cdots+\dfrac{1}{2^n}=1-\dfrac{1}{2^n}$, which tend to $1$, so the series converges and $S$ is a genuine number. For convergent series, multiplying termwise by a constant and adding or subtracting termwise are allowed, because they are just limits of the same operations on partial sums: $2S_n=1+S_{n-1}$ for every $n$, and letting $n\to\infty$ gives $2S=1+S$. So $S=1$ is correct (and agrees with $S_n\to1$).

\textbf{Second series.} The partial sums are $T_n=1,0,1,0,\dots$, which do not approach any limit, so $T$ is not a number at all; the series diverges. The step $T=1-T$ treats $T$ as a number and rearranges it, which is meaningless. Assuming that a divergent series has a value leads to contradictions: grouping as $(1-1)+(1-1)+\cdots$ would give $T=0$, while $1+(-1+1)+(-1+1)+\cdots$ would give $T=1$.

So the manipulation is justified for convergent series (where linearity of limits applies), not for divergent ones. (The value $\dfrac{1}{2}$ does have a meaning in more general summation methods, such as the average of the partial sums, but not as an ordinary sum.)
